\documentclass[answers]{exam}

% Version compilable
%\usepackage{../../../mypackages}
%\usepackage{../../../macros}

% Version pour execution du script
\usepackage{../../mypackages}
\usepackage{../../macros}

% Pour le tableau
\usepackage{array}
\usepackage{longtable}
\usepackage{multirow}
\usepackage[table]{xcolor}

\title{Notation DM Geogebra}
\author{Nicolas Bancel}
\date{6 Décembre 2025}

% Mise en forme des solutions en bleu
\SolutionEmphasis{\color{blue}}
\renewcommand{\solutiontitle}{\noindent}

\definecolor{lavande}{RGB}{180,190,255}

% Colonne p{} alignée à gauche et qui wrappe
\newcolumntype{L}[1]{>{\raggedright\arraybackslash}p{#1}}

\begin{document}

\textbf{Collège Lycée Suger}
\hfill
\textbf{Mathématiques} \\

\textbf{Année 2025-2026}
\hfill
\textbf{Classe de 6ème} \par

{\let\newpage\relax\maketitle}

\subsection*{Barème de notation du devoir maison}

\renewcommand{\arraystretch}{1.7}

{
  %\setlength{\LTleft}{-2.3cm}  % déborde de 1.5 cm dans la marge de gauche
  \setlength{\LTright}{0pt}    % rien de plus à droite

\centering
\begin{longtable}{|L{0.3\textwidth}|L{0.7\textwidth}|c|}
\hline
\rowcolor{lavande}
\textbf{Critère évalué} & \textbf{Description} & \textbf{Points} \\[2pt]
\hline

% ---------------- Travail mathématique et géométrique ----------------
  Exercice d'application n°3 (fichier GeoGebra)
  & L'exercice d'application n°3 est réalisé complètement. 
  La construction (angles, longueurs, nom des points) est respectée.
  Il n'y a pas de point en trop, qui ne sert à rien dans la figure.
  & 4 \\
\hline

  Figure sur papier (constellation)
  & La figure de la constellation est réalisée au crayon à papier, avec les instruments de géométrie, sans tracés à main levée.
  Les points, segments et éventuelles légendes sont clairement indiqués et lisibles.
  La figure a un titre : vous indiquez clairement quelle constellation vous avez choisi.
  & 3 \\
\hline

  Figure GeoGebra (constellation)
  & La figure GeoGebra respecte les consignes : construction propre, segment ou angles correctement tracés, et nommés.
  Il n'y pas de point en trop, qui ne sert à rien dans la figure et qui n'en fait pas partie. 
  Le fichier est enregistré et nommé correctement (vous pouvez préciser la constellation que vous avez choisie)
  & 3 \\[6pt]
\hline  

% ---------------- Organisation et outils numériques ----------------

Utilisation de Gmail + Autonomie + Bonne volonté
& Le devoir est envoyé à la bonne adresse, avec un objet de mail clair ("Le titre de votre email doit être "VotrePrénom VotreNom 6ème - Devoir maison Geogebra", par exemple "Nicolas Bancel 6ème - Devoir maison Geogebra").
  Les 3 pièces jointes (2 fichiers GeoGebra + 1 photo de la figure papier) sont bien attachées. 
  Capacité à utiliser Internet et se débrouiller seul pour trouver les réponses à ses questions.
  Ou si un email sur EcoleDirecte est envoyé pour poser une question, la question est claire et précise.
& 5 \\ [6pt]
\hline

% ---------------- Gestion du temps et anticipation ----------------

  Anticipation du travail
  & Sur 5 points :  
  \begin{compactitem}
    \item 5 pts : travail commencé et/ou questions posées avant la date limite, prise en compte des feedbacks, devoir envoyé en avance.
    \item 2.5 pts : devoir envoyé dans les temps mais questions posées la veille de la deadline ou le jour même. Ou non prise en compte des retours que j'ai pu faire.
    \item 0 pt : devoir envoyé le jour même mais ou pendant l'heure de cours. 
  \end{compactitem}
  \textbf{Malus de retard :} si le devoir est rendu en retard, on applique un \textbf{–5} sur la note finale du devoir, en plus des points d’anticipation.
  & 5 \\[6pt]
\hline

% ---------------- Total ----------------
\rowcolor{green!35}
 \textbf{Total} & \textbf{Note maximale possible (hors malus de retard)} & \textbf{20} \\
\hline

\end{longtable}
}

\end{document}
